\documentclass{article}
\usepackage{graphicx} % Required for inserting images
\usepackage[
backend=biber,
style=numeric,
sorting=ynt
]{biblatex}
\addbibresource{references.bib}
\usepackage{amsmath}
\usepackage{amssymb}

\title{MSE342 - Project Proposal}
\author{Sondre Rogde}
\date{May 2026}

\begin{document}

\maketitle

\section{Executive Summary}
Generate financial time series using diffusion models with a topological loss term. 

\section{Related Work}
The seminal work of Gidea \& Katz in \cite{Gidea_2018} found that topological features can be predictive of imminent market crashes. They find that $L^p$-norms of the persistence landscape of point clouds over sliding windows increase a lot for up to a year before a market crash. They use sliding windows of time series data as point clouds and look at how the $L^p$ norm of the persistence landscape of these evolve over time. 

Persistence landscape proposed in \cite{Bubenik_2015} is a topological summary designed be compatible with machine learning. Using the persistence landscape allows us to incorporate a topological loss term when doing gradient descent for the diffusion model.

To make use of the topological features it makes sense to use more than just price data. Diffusion models have also been used with richer data than just price data in \cite{Takahashi_2025}. 


\section{Mathematical Framing}

\subsection*{Sliding Window Point Clouds}
Following \cite{Gidea_2018}, given a multivariate time series $x(t_n) = (x_n^1, \dots, x_n^d) \in \mathbb{R}^d$ for $n = 1, \dots, T$, we define a sequence of point clouds via sliding windows of width $w$:
\[
    X_n = \bigl(x(t_n),\, x(t_{n+1}),\, \dots,\, x(t_{n+w-1})\bigr) \subset \mathbb{R}^d, \quad n = 1, \dots, T - w + 1.
\]

\subsection*{Persistence Diagrams and Landscapes}
For each point cloud $X_n$ we compute the Vietoris--Rips filtration $\{VR_\varepsilon(X_n)\}_{\varepsilon \geq 0}$ and extract persistence diagrams $\mathrm{Dgm}_k(X_n) = \{(b_i, d_i)\}_i$ for homological dimensions $k = 0, 1$, capturing connected components and loops respectively.

Following \cite{Bubenik_2015}, the $j$-th persistence landscape $\lambda_j : \mathbb{R} \to \mathbb{R}$ is defined from $\mathrm{Dgm}_k(X_n)$ as the $j$-th largest value of the tent functions
\[
    f_i(t) = \max\!\bigl(0,\, \min(t - b_i,\, d_i - t)\bigr),
\]
giving a sequence $\lambda_1 \geq \lambda_2 \geq \cdots$ in $L^2(\mathbb{R})$. We write $\Lambda_k(X_n) = (\lambda_1^{(k)}, \lambda_2^{(k)}, \dots)$ for the full landscape in dimension $k$. The map from point cloud to persistence landscape is differentiable almost everywhere \cite{Carriere_2021}, enabling gradient-based optimization.

\subsection*{Score-Based Diffusion Model}
We follow the SDE framework of \cite{Song_2021}. A data sample $x_0 \sim p_{\mathrm{data}}$ is corrupted by the forward process
\[
    dX_t = f(X_t, t)\,dt + g(t)\,dW_t, \quad t \in [0, T],
\]
where $W_t$ is a standard Wiener process. A neural network $s_\theta(x, t)$ approximates the score $\nabla_x \log p_t(x)$ by minimizing the denoising score matching objective
\[
    \mathcal{L}_{\mathrm{DSM}}(\theta) = \mathbb{E}_{t,\,x_0,\,x_t}\!\left[ w(t)\,\bigl\|s_\theta(x_t, t) - \nabla_{x_t} \log p(x_t \mid x_0)\bigr\|^2 \right],
\]
where $w(t) > 0$ is a weighting function. New samples are generated by integrating the reverse-time SDE.

\subsection*{Topological Loss}
Let $\hat{x}_0 \sim p_\theta$ denote a generated sample with corresponding sliding-window clouds $\hat{X}_n$. We define the topological loss as the $L^2$ discrepancy between the empirical mean persistence landscapes of real and generated data:
\[
    \mathcal{L}_{\mathrm{topo}}(\theta) = \sum_{k=0}^{1} \left\| \frac{1}{N}\sum_{n=1}^N \Lambda_k(X_n) - \frac{1}{N}\sum_{n=1}^N \Lambda_k(\hat{X}_n) \right\|_{L^2}^2.
\]
The full training objective is
\[
    \mathcal{L}(\theta) = \mathcal{L}_{\mathrm{DSM}}(\theta) + \alpha\,\mathcal{L}_{\mathrm{topo}}(\theta),
\]
where $\alpha > 0$ balances the two terms. Because the persistence landscape computation is differentiable, gradients of $\mathcal{L}_{\mathrm{topo}}$ flow back through the generator.

\section{Data}
For this project we need at least the price of a stock. However, more data such as 

A possible dataset could be $5$ liquid stocks (i.e. AAPL, AMZN, MSFT, JPM, GOOGL) featuring prices, volume, volatility, bid-ask spread, price earnings, maybe other accounting-time series hybrids. We want 10-15 years of data, involving at least one crash. 

\section{Risks \& Mitigation}
\begin{itemize}
    \item If it's too easy or we finish faster than expected we can use the synthetic data for downstream tasks such as optimal stopping for American options. 
    \item We don't get good results on daily data. Can look for intraday data.
\end{itemize}

\section{Other loosely connected things}
\begin{itemize}
    \item Can use Wasserstein distance as distance metric between persistence diagrams. This allows us to use more of the material from class.
\end{itemize}

\printbibliography


\end{document}
